\PassOptionsToPackage{hyperfootnotes=false}{hyperref}
\documentclass{beamer}
\usefonttheme{professionalfonts}
%\usetheme{Rochester}
%\usetheme{AnnArbor}
\usetheme{Szeged}
\useoutertheme{infolines}
\usecolortheme{dolphin}
\usepackage{graphicx}
% \usepackage[english]{babel}
\usepackage{multicol}
\usepackage{color}
\usepackage{xcolor}
%\usepackage{smartdiagram}


%\definecolor{carmine}{rgb}{0.59, 0.0, 0.09}
\usepackage{amsmath} 

\usepackage{graphicx}
\usepackage{tikz}
\usetikzlibrary{shapes.callouts}
\usepackage{multirow}

\usepackage[utf8]{inputenc}

\usepackage{enumerate}
\usepackage[shortlabels]{enumitem}
\usepackage{amsmath,amsthm,amssymb,amsfonts,mathrsfs,latexsym,stmaryrd}
\usepackage{tcolorbox}

\usepackage{color} 
\usepackage{amssymb, amsmath, amsbsy}
\usepackage{adjustbox,lipsum}
\usepackage{placeins}

\usepackage{textpos}
\usepackage{wrapfig}
\usepackage[spanish, es-tabla]{babel}

\usepackage{caption}
\setlength{\parindent}{0pt}

\usepackage{url}
\usepackage{float}
\usepackage{tikz}
\usetikzlibrary{patterns}


\title[ICS-162]{Econometría ICS-294\\ \small{Clase 15: Test de Chow y variables dependientes binarias }}

\author{Francisco Alfaro Medina}
\institute[USM]
{Universidad Técnica Federico Santa María\\ % Your institution for the title page
	\medskip
	\textit{} % Your email address
	
}
\logo{

\includegraphics[width=1cm]{Logo.png}
}
\date{}
\newcommand\Fontvi{\fontsize{9}{7.2}\selectfont}
\begin{document}
	
\begin{frame}
	\titlepage % Print the title page as the first slide


\end{frame}
%%%%%%%%%%%%%%%%%%%%%%%%%%%%%%%%%%%%%%%%%%%%%%%%%%%%%%%%%%%%%%%%%%%%%%%%%%%

\begin{frame}[shrink=10]{ Diferencias entre grupos}
\begin{itemize}
    \item Suponga el siguiente modelo:
    \[
    y = \beta_0 + \beta_1 x_1 + \beta_2 x_2 + \beta_3 x_3 + u
    \]
    y que se puede clasificar las \(n\) observaciones en \(G\) grupos distintos.
\end{itemize}

\begin{itemize}
    \item Se quiere testear que el modelo es el mismo para todos los grupos, es decir, los coeficientes son los mismos.
    \item En series de tiempo esto equivale a testear un cambio de régimen o cambio “estructural”.
\end{itemize}

\begin{itemize}
    \item Suponga \(G = 2\). Para realizar el test planteamos el modelo sin restringir:
    \[
    y = \beta_0 + \beta_1 x_1 + \beta_2 x_2 + \beta_3 x_3 + \gamma_0 D + \gamma_1 x_1 D + \gamma_2 x_2 D + \gamma_3 x_3 D + u
    \]
    donde \(D\) es una variable dummy para observaciones de un mismo grupo (e.g: mujer).
\end{itemize}

\begin{itemize}
    \item La hipótesis nula correspondiente es:
    \[
    H_0 : \gamma_0 = \gamma_1 = \gamma_2 = \gamma_3 = 0
    \]
    \item Y se puede hacer un test F.
\end{itemize}
\end{frame}
\begin{frame}[shrink=10]{Test de Chow: diferencias entre grupos}
\begin{itemize}
    \item Otra forma de verificar esta hipótesis es usando el Test de Chow que plantea hacer una regresión para cada grupo.
    \item Estimar la ecuación:
    \[
    y = \beta_{0}^g + \beta_{1}^g x_1 + \beta_{2}^g x_2 + \beta_{3}^g x_3 + u
    \]
    para cada subpoblación \(g = 1, \ldots, G\).
\end{itemize}

\begin{itemize}
    \item Obtener la \(SSR_g\) correspondiente a la regresión para dicho subgrupo, y construir \(SSR_{ur} = \sum_{g=1}^G SSR_g\).
    \item Estimar la ecuación restringida:
    \[
    y = \beta_0 + \beta_1 x_1 + \beta_2 x_2 + \beta_3 x_3 + u
    \]
    usando todas las observaciones y obtener su \(SSR_r\).
\end{itemize}

\begin{itemize}
    \item Construir el estadístico del Test de Chow:
    \[
    F = \frac{(SSR_r - SSR_{ur}) / ((G - 1)(k + 1))}{SSR_{ur} / (n - G(k + 1))} \sim F((G - 1)(k + 1), n - G(k + 1))
    \]
\end{itemize}
\end{frame}

\begin{frame}{Ejemplo test de Chow }
% \begin{figure}
%     \centering
%     \includegraphics[width=0.8\linewidth]{sem11//images/chow.png}
% \end{figure}


    \[
\log(wage) = \beta_0 + \beta_1 educ + \beta_2 experiencia + \beta_3 ingreso\_familiar + u
\]
{\footnotesize
\begin{itemize}
    \item Queremos saber si el modelo es igual para dos grupos ($G=2$): mujeres (g=1) y hombres (g=2). La hipótesis nula del test de Chow es que los modelos son iguales.
    
  
        \item Hacemos la regresión para cada subgrupo:
        \begin{itemize}
            \item Solo considerando las 90 mujeres de la muestra, obtenemos             $SSR_1 = 19$
            \item Solo considerando los 276 hombres de la muestra, obtenemos 
            $SSR_2 = 59$
        \end{itemize}
        
        \item Sumamos estos dos números para obtener:         $SSR_{ur} = 78$
        
        \item Regresión restringida: hacemos la regresión para toda la muestra de 366 personas (asumiendo que no hay diferencias entre hombres y mujeres). Obtenemos    $SSR_r = 85$
        
        \item Construimos el F-estadístico de Chow:
        \[
        F = \frac{85 - 78}{78} \cdot \frac{366 - 2(3 + 1)}{(2 - 1)(3 + 1)} = 8,03
        \]
        
        \item Comparamos al 5\% con $F(4, 358) = 2.39$. Se rechaza la hipótesis nula, es decir, hay diferencias entre los grupos.
  
\end{itemize}
}
 \end{frame}

\begin{frame}{Ejercicio}

{\small Una ecuación que explica el sueldo de los presidentes de consejos de administraci\'on es:


{\tiny
\begin{align*}
\log(\widehat{salary}) &= 4.59 + 0.257 \log(sales) + 0.011 \, roe + 0.158 \, finance 
 + 0.181 \, consprod - 0.283 \, service \\
&\quad (0.30) \quad \quad \quad (0.032) \quad  \quad \quad  \quad \quad  (0.004) \quad \quad \quad \quad \quad (0.089) \quad\quad \quad \quad \quad  (0.085) \quad \quad \quad \quad  (0.099) \\
n &= 209, \quad R^2 = 0.357
\end{align*}
}
% \begin{figure}
%     \centering
%     \includegraphics[width=0.65\linewidth]{sem11//images/ejer.png}
% \end{figure}

donde finance, consprod y
service son variables binarias que corresponden a las industrias financiera, de productos de consumo y de servicios. La industria que se ha omitido es la del transporte.}
\begin{itemize}
\item {\small Calcule la diferencia porcentual aproximada entre los sueldos estimados de las industrias
de servicios y de transporte, manteniendo sales y roe constantes. Es esta diferencia estadísticamente significativa?}
\item {\small Cu\'al es la diferencia porcentual aproximada entre los sueldos estimados de la industria
de productos de consumo y en el sector financiero? C\'omo podría probar
si esta diferencia es estad\'isticamente significativa?}

\end{itemize}
    
\end{frame}


\begin{frame}{Variables dependientes binarias}

    Muchos problemas de la vida real consisten en explicar una variable dummy. 
\begin{itemize}
    \item ¿Qu\'e explica si un alumno va a un colegio p\'ublico o privado.
\item ¿Qui\'enes son las personas que se jubilan anticipamente?
\item  ¿Qu\'e tipos de personas tienen mas probabilidad de chocar el auto?
\end{itemize}
 
\end{frame}

\begin{frame}{Modelo de probabilidad lineal (MPL)}

    % \begin{figure}
    %     \centering
    %     \includegraphics[width=0.9\linewidth]{sem11/images/MPL.png}
    % \end{figure}
{\footnotesize
    \begin{itemize}
    \item Considere el caso simple cuando $y$ tiene sólo dos categorías, definidas como 0 (fracaso) y 1 (éxito).
    
    \begin{align*}
    y &= \beta_0 + \beta_1 x_1 + \beta_2 x_2 + \dots + \beta_k x_k + u \\
    \mathbb{E}[y \mid x] &= \beta_0 + \beta_1 x_1 + \beta_2 x_2 + \dots + \beta_k x_k
    \end{align*}
    
    \item En este caso se cumple además que:
    
    \[
    \mathbb{E}[y \mid x] = 1 \cdot \mathbb{P}[y = 1 \mid x] + 0 \cdot \mathbb{P}[y = 0 \mid x] = \mathbb{P}[y = 1 \mid x]
    \]
    
    \item Luego, la regresión es igual a:
    
    \[
    \mathbb{P}[y = 1 \mid x] = \beta_0 + \beta_1 x_1 + \beta_2 x_2 + \dots + \beta_k x_k
    \]
    
    \item Los modelos con variable dependiente binaria son llamados \underline{modelos de probabilidad lineal}.
    
    \item $\beta_j \quad (j\neq 0)$ mide el cambio en la probabilidad de éxito debido a un cambio en una unidad de $x_j$.
    $\beta_0$ mide la probabilidad de éxito cuando $x_1 = x_2 = \dots = x_k = 0$.
    
\end{itemize}}
\end{frame}

\begin{frame}{El MPL presenta dos problemas}

\begin{enumerate}
    \item  Se pueden predecir individuos que tengan probabilidades de \'exito negativas o
mayores que uno.
Si nos interesa el efecto en el promedio, este problema no es muy grave.
Si nos interesa la predicción o efecto para individuos extremos, los modelos
no-lineales Logit y Probit son mejores para este tipo de casos (no los vamos a ver
en el curso).
\item No se cumple el supuesto de homocedasticidad.
\end{enumerate}

\end{frame}

\begin{frame}{Ejemplo}

Considere la siguiente regresi\'on, donde la variable dependiente {\color{blue}Accidente}
toma valor 1 si la persona tuvo al menos un accidente de tr\'ansito y 0 si no.
$$\widehat{Accidente}_i = 0.30 -0.05educ_i -0.0001ingreso_i $$
\begin{itemize}
    \item C\'omo interpretar el $\hat{\beta}_0$=0,30? 

La probabilidad de accidente es 30\% para
individuos con 0 años de educaci\'on y 0 ingreso.
    \item C\'omo interpretar el coeficiente de educ? 

Un aumento de la educaci\'on de la
persona en un año reduce la probabilidad de accidente en 0,05.
    \item  Cu\'al es la probabilidad de accidente de una persona cuyo ingreso es 2000
d\'olares y tiene 20 años de educaci\'on? 


\item Notar que si ponen valores altos de educ e ingreso se obtiene una
probabilidad negativa!
\end{itemize}
    
\end{frame}


\end{document}
